% GENERATED FILE. Edit canonical lessons, then run npm run generate:books.
% canonical-source-hash: fnv1a64:97e8c1c1702c386b
% canonical-lessons: GE-C203-portemonnaie, GE-C203-groesse, GE-C203-stoff, GE-C203-wolle, GE-C203-seide

\chapter{A purse, A size, Material, Wool, Silk}
\label{ch:ge-a-purse-a-size-material-wool-silk}
\hlchaptermodality{203}

\begin{chapteropening}
\textbf{By the end of this chapter:} \emph{I can say a purse, a size, material, wool and silk.}

Complete the last lesson of chapter 203: I can say a purse, a size, material, wool and silk.
\end{chapteropening}

\section[das Portemonnaie]{das Portemonnaie — a purse, a wallet}

\label{lesson:GE-C203-portemonnaie}



\begin{quote}
Before the new one: say the German for a brush, then the German for a suit.
\end{quote}

\subsection*{You'll want to know: das Portemonnaie}
\textbf{das Portemonnaie} — \textquotedblleft{}a purse, a wallet\textquotedblright{}.

\begin{cousinweb}
\textbf{Portemonnaie} $\leftarrow$ French \emph{porte-monnaie}, \textquotedblleft{}carry-money\textquotedblright{} (EtymDB 2.1) — borrowed.
\end{cousinweb}

\begin{grammarlens}[title={What you've built}]
One more word for this chapter.
\end{grammarlens}

\subsection*{Your turn}
\begin{itemize}
\raggedright
  \item \emph{Say it:} \emph{das Portemonnaie}
  \item \emph{Say it:} \emph{das Portemonnaie}, once more
  \item \emph{Say it:} say \emph{das Portemonnaie}
  \item \emph{Recall:} say the German for a brush, then the German for a suit, then say \emph{das Portemonnaie} again
\end{itemize}

\begin{tcolorbox}[breakable,colback=teal!4,colframe=teal!35!black,title={Before you move on}]
What does das Portemonnaie mean? (\textquotedblleft{}A purse, a wallet\textquotedblright{}.) Say it once more.
\end{tcolorbox}
\section[die Größe]{die Größe — a size}

\label{lesson:GE-C203-groesse}



\begin{quote}
Before the new one: say the German for a suit, then the German for a purse.
\end{quote}

\subsection*{You'll want to know: die Größe}
\textbf{die Größe} — \textquotedblleft{}a size\textquotedblright{}.

\begin{cousinweb}
\textbf{Größe} $\leftarrow$ Old High German \emph{grōzī} $\leftarrow$ Proto-Germanic \emph{*grautį̄}, \textquotedblleft{}greatness, size\textquotedblright{} — inherited, the noun of the word for \textquotedblleft{}big, great\textquotedblright{}, a cousin of English \emph{great}. (EtymDB 2.1)
\end{cousinweb}

\begin{grammarlens}[title={What you've built}]
One more word for this chapter.
\end{grammarlens}

\subsection*{Your turn}
\begin{itemize}
\raggedright
  \item \emph{Say it:} \emph{die Größe}
  \item \emph{Say it:} \emph{die Größe}, once more
  \item \emph{Say it:} say \emph{die Größe}
  \item \emph{Recall:} say the German for a suit, then the German for a purse, then say \emph{die Größe} again
\end{itemize}

\begin{tcolorbox}[breakable,colback=teal!4,colframe=teal!35!black,title={Before you move on}]
What does die Größe mean? (\textquotedblleft{}A size\textquotedblright{}.) Say it once more.
\end{tcolorbox}
\section[der Stoff]{der Stoff — material, fabric}

\label{lesson:GE-C203-stoff}



\begin{quote}
Before the new one: say the German for a purse, then the German for a size.
\end{quote}

\subsection*{You'll want to know: der Stoff}
\textbf{der Stoff} — \textquotedblleft{}material, fabric\textquotedblright{}.

\begin{cousinweb}
\textbf{Stoff} came in from Old French \emph{estoffe}, \textquotedblleft{}material\textquotedblright{}; EtymDB 2.1 lists it beside English \emph{stuff}, the same French word.
\end{cousinweb}

\begin{grammarlens}[title={What you've built}]
One more word for this chapter.
\end{grammarlens}

\subsection*{Your turn}
\begin{itemize}
\raggedright
  \item \emph{Say it:} \emph{der Stoff}
  \item \emph{Say it:} \emph{der Stoff}, once more
  \item \emph{Say it:} say \emph{der Stoff}
  \item \emph{Recall:} say the German for a purse, then the German for a size, then say \emph{der Stoff} again
\end{itemize}

\begin{tcolorbox}[breakable,colback=teal!4,colframe=teal!35!black,title={Before you move on}]
What does der Stoff mean? (\textquotedblleft{}Material, fabric\textquotedblright{}.) Say it once more.
\end{tcolorbox}
\section[die Wolle]{die Wolle — wool}

\label{lesson:GE-C203-wolle}



\begin{quote}
Before the new one: say the German for a size, then the German for material.
\end{quote}

\subsection*{You'll want to know: die Wolle}
\textbf{die Wolle} — \textquotedblleft{}wool\textquotedblright{}.

\begin{cousinweb}
\textbf{Wolle} $\leftarrow$ Old High German \emph{wolla} — inherited, the same word as English \emph{wool}. EtymDB 2.1 has the German word but not its source.
\end{cousinweb}

\begin{grammarlens}[title={What you've built}]
One more word for this chapter.
\end{grammarlens}

\subsection*{Your turn}
\begin{itemize}
\raggedright
  \item \emph{Say it:} \emph{die Wolle}
  \item \emph{Say it:} \emph{die Wolle}, once more
  \item \emph{Say it:} say \emph{die Wolle}
  \item \emph{Recall:} say the German for a size, then the German for material, then say \emph{die Wolle} again
\end{itemize}

\begin{tcolorbox}[breakable,colback=teal!4,colframe=teal!35!black,title={Before you move on}]
What does die Wolle mean? (\textquotedblleft{}Wool\textquotedblright{}.) Say it once more.
\end{tcolorbox}
\section[die Seide]{die Seide — silk}

\label{lesson:GE-C203-seide}



\begin{quote}
Before the new one: say the German for material, then the German for wool.
\end{quote}

\subsection*{You'll want to know: die Seide}
\textbf{die Seide} — \textquotedblleft{}silk\textquotedblright{}.

\begin{cousinweb}
\textbf{Seide} $\leftarrow$ Middle High German \emph{sīde} $\leftarrow$ Old High German \emph{sīda}, \textquotedblleft{}silk\textquotedblright{} (EtymDB 2.1), from Medieval Latin \emph{sēta}, \textquotedblleft{}bristle, silk\textquotedblright{} — borrowed.
\end{cousinweb}

\begin{grammarlens}[title={What you've built}]
One more word for this chapter.
\end{grammarlens}

\subsection*{Your turn}
\begin{itemize}
\raggedright
  \item \emph{Say it:} \emph{die Seide}
  \item \emph{Say it:} \emph{die Seide}, once more
  \item \emph{Say it:} say \emph{die Seide}
  \item \emph{Recall:} say the German for material, then the German for wool, then say \emph{die Seide} again
\end{itemize}

\begin{tcolorbox}[breakable,colback=teal!4,colframe=teal!35!black,title={Before you move on}]
What does die Seide mean? (\textquotedblleft{}Silk\textquotedblright{}.) Say it once more.
\end{tcolorbox}
